% TikZ 图：高中实验常用基本方法
% Day 28：实验专题
\begin{tikzpicture}[scale=0.95,>=Latex]
  \tikzset{
    card/.style={rounded corners=5pt,thick,fill=gray!8,text width=4.05cm,minimum height=1.55cm,align=center,font=\small},
    arrow/.style={->,thick,blue}
  }

  \node[card] (control) at (0,1.15) {
    \textbf{控制变量法}\\[3pt]
    只改变一个量\\
    研究 $a$ 与 $F$、$m$ 的关系
  };
  \node[card] (convert) at (5.2,1.15) {
    \textbf{转换放大法}\\[3pt]
    把不易测的量转为易测量\\
    纸带、光杠杆、传感器
  };
  \node[card] (equiv) at (0,-1.15) {
    \textbf{等效替代法}\\[3pt]
    标准对象与待测对象同效\\
    电阻箱替代未知电阻
  };
  \node[card] (compare) at (5.2,-1.15) {
    \textbf{比较法}\\[3pt]
    与标准量或理论值比较\\
    电桥平衡、机械能验证
  };

\end{tikzpicture}
